\section{Conclusion}\label{sec:conclusion}

Declared and assessed \vat{} debt in S\~ao Paulo differ in size, concentration, payment and recovery by enough that compliance studies should not pool them by default.
Assessed debt is a small number of large, older certificates and most of the value.
Declared debt is a large number of small certificates, usually several on the same firm, and a higher recovery relative to the stock that remains.
The largest 1~percent of certificates hold three quarters of the pooled stock because that tail is assessed debt, not because each segment is concentrated to that degree.
For the attorney office, the practical content is a specialised list of large assessed certificates on one side, and a multi-certificate declared book on the other.

These descriptions come from one cross-section of enrolled debt, and the firm register is observed after the outcomes, so a closed or unfit status can be a consequence of the debt rather than a prior trait.
Certificates of the same firm are also dependent, which is why the firm-level effect sizes sit next to the certificate-level ones.
A useful next step is to follow the same certificates through payment and instalment, and to ask whether a review concentrated on the largest certificates changes the amount recovered.
Whether enrolled VAT stocks in other administrations have the same shape is an empirical question for those stocks, and the S\~ao Paulo evidence says it is worth asking before an enrolled portfolio is treated as one residual.
